\section{Polynomial moments and decay of irrelevant insertions}
\label{sec:polynomial-moment-attenuation}
The subtraction in Section~\ref{sec:local-splitting} removes the
low-order holonomy tests of an insertion. We now explain why this
improves its coarse-to-fine estimate. The argument has two parts:
a qualitative implication from holonomy tests to polynomial moments,
and a quantitative decomposition of the actual interpolating field.
Both use the same centered lattice lift as the local subtraction.

\subsection{From holonomy tests to polynomial moments}
Let $X$ be a nonwrapping support about $o$ on a torus of side at
least two. Let $f$ be local on $X$, gauge invariant, and $C^d$ in
logarithmic coordinates near the identity. Suppose that every
holonomy test through degree $d$ vanishes:
\begin{equation}\label{eq:moment-test-null}
 \left.\frac{d^j}{dt^j} f(H_o(J,t))\right|_{t=0}=0,
                   \qquad 0\le j\le d,\quad\hbox{every }J.
\end{equation}
Tests with a larger available cutoff are allowed; only degrees
through $d$ enter the argument. For $a\ge1$ let $\mathcal M_a$
be the space of link fields screened to the links with both
endpoints in $X$, whose value on a link in direction $i$ is a
polynomial of degree at most $a-1$ in
$x=\operatorname{lift}(e_--o)$. We claim that
\begin{equation}\label{eq:polynomial-moment-null}
 D^m(f\circ\exp)(0)[v_1,\ldots,v_m]=0
 \quad\text{if }v_i|_X\in\mathcal M_{a_i},\quad
 a_i\ge1,\quad\sum_i a_i\le d.
\end{equation}
The restriction means screening the link field, not restricting
its support before applying the coarse-to-fine map.

Here are the details of the change of coordinates behind this
claim. Introduce a site gauge parameter $q$ and the finite list
$J$ of radial connection coefficients of weights $2,3,\ldots$.
For each link set
\[
 \Psi(q,J)_e=
 \log\bigl(e^{q(e_-)}H_o(J,1)_e e^{-q(e_+)}\bigr)
\]
on the links tested by $X$, and set it to zero elsewhere. This
is an analytic germ with $\Psi(0)=0$. The endpoint map is analytic
by the integral-equation argument in Section~\ref{sec:local-splitting}.
Locality and gauge invariance give, as germs,
\begin{equation}\label{eq:moment-gauge-germ}
            (f\circ\exp)\circ\Psi(q,J)=f(H_o(J,1)).
\end{equation}
The identity uses only tested links and a finite common chart.

Give a direction $(q,J)$ filtration degree $a$ when $q$ is
polynomial of degree at most $a$ on $X$ and the components of
$J$ of weight greater than $a$ vanish. Denote this space by
$\mathcal K_a$. The following two properties hold:
\begin{align}
 D\Psi(0)\mathcal K_a&=\mathcal M_a,
       &&1\le a\le d,\label{eq:moment-surjection}\\
 D^r\Psi(0)[u_1,\ldots,u_r]&\in\mathcal M_{\sum_i a_i},
       &&u_i\in\mathcal K_{a_i},\ r\ge1.
       \label{eq:moment-filtration}
\end{align}
We prove them, since the nonlinear coordinate change matters.

For the first property the linearized link is
\[
 D\Psi(0)(q,J)_{(x,i)}
       =q(x)-q(x+e_i)+\int_0^1 A_J(x+se_i)(e_i)\,ds.
\]
Nonwrapping and side at least two identify the centered lift at
the endpoint with $x+e_i$. Every constant one-form is the gradient
of a linear site polynomial, so it occurs with $J=0$. For a
homogeneous polynomial one-form $Q$ of degree $k\ge1$, put
\[
 r(x)=\frac1{k+1}\sum_i x_iQ_i(x),\qquad K_i=Q_i-\partial_i r.
\]
Euler's identity gives $\sum_i x_iK_i(x)=0$. Thus $K$ is precisely
an admissible radial connection coefficient of weight $k+1$.
Taking $q=-r$ and that coefficient gives the link average of $Q$.
The difference between this average and $Q_i(x)$ has degree at
most $k-1$: expand $Q_i(x+se_i)$ as a polynomial in $s$, and
observe that its constant term is $Q_i(x)$. Induction on the
degree supplies the difference by lower-weight parameters.
Splitting a general polynomial into homogeneous parts proves
\eqref{eq:moment-surjection}.

For \eqref{eq:moment-filtration}, write the link germ in variables
\[
     q(x),\qquad q(x+e_i)-q(x),\qquad A_J(x+se_i).
\]
For a direction of degree $a$, the first expression has degree
at most $a$ in $x$, and each of the other two has degree at most
$a-1$. When both of the latter expressions vanish, the link germ
is identically zero, since $e^{q(x)}1e^{-q(x)}=1$. Expand any
multilinear derivative according to these three components in
each argument. The term using only the first components vanishes.
Every remaining term contains a component with a one-degree
loss; its degree is at most $\sum_i a_i-1$. Screening completes
\eqref{eq:moment-filtration}.

Equation~\eqref{eq:moment-test-null} implies that all mixed
derivatives of the right side of \eqref{eq:moment-gauge-germ}
with total filtration weight at most $d$ vanish. To verify this
precisely, take its Taylor polynomial through ordinary degree
$d$ in the finite coordinates $J$. Substitute the weighted
scaling $J_a\mapsto t^aJ_a$. The coefficient of every $t^j$,
$j\le d$, is zero for all $J$, by \eqref{eq:moment-test-null}.
Each such coefficient is a polynomial, hence all its monomial
coefficients vanish. Polarization gives the asserted mixed
derivatives. Gauge-coordinate derivatives vanish because the
right side is independent of $q$.

Apply the chain rule to the left side of
\eqref{eq:moment-gauge-germ}. Induct on the number $m$ of
arguments in \eqref{eq:polynomial-moment-null}, lifting each
$v_i$ to $u_i\in\mathcal K_{a_i}$ by
\eqref{eq:moment-surjection}. The partition into singletons
contributes exactly the desired derivative of $f\circ\exp$.
Every other partition has fewer than $m$ blocks; by
\eqref{eq:moment-filtration}, its arguments have the block
weights, whose sum is still $\sum_i a_i$. The induction hypothesis
kills all these terms. The whole derivative of the composite
is zero, so the singleton term is zero too. The case $m=0$ is
$f(1)=0$. Finally, locality makes $f\circ\exp$ equal to its
composition with the linear screening map. Differentiating that
identity allows unscreened directions whose restrictions have
the stated polynomial form. This proves
\eqref{eq:polynomial-moment-null}.

\subsection{A discrete Taylor decomposition on the cover}
Write $|x|_1=\sum_{i=1}^4|x_i|$ for $x\in\mathbb Z^4$ and
$\Delta_i v(x)=v(x+e_i)-v(x)$. For a word $w$ in the four
axes let $\Delta_w$ be the corresponding iterated difference.
Suppose $d,p\in\mathbb N$, $L\ge1$, $p\ge d$, $M\ge0$,
$\theta>0$, and a colour-valued cover field $v$
satisfies
\begin{equation}\label{eq:moment-drift}
 \|\Delta_wv(x)\|\le M\theta L^{-1-|w|}
                       (1+|x|_1/L)^p,
                  \qquad |w|\le d.
\end{equation}
There are polynomial pieces $P_j$, $0\le j<d$, and a remainder
$R_d$ such that
\begin{equation}\label{eq:moment-newton}
 v(x)=\sum_{j<d}P_j(x)+R_d(x),\qquad
 \begin{cases}
 \deg P_j\le j,\\
 \|P_j(x)\|\le C_{d,p}M\theta L^{-j-1}(1+|x|_1)^p,\\
 \|R_d(x)\|\le C_{d,p}M\theta L^{-d-1}(1+|x|_1)^p.
 \end{cases}
\end{equation}
The constant precedes the choice of $L$ and of the period.

For completeness use the total-degree Newton polynomials
\[
 P_j(x)=\sum_{|\alpha|=j}\binom{x}{\alpha}\Delta^\alpha v(0),
 \qquad
 \binom{x}{\alpha}=\prod_{i=1}^4
       \frac{x_i(x_i-1)\cdots(x_i-\alpha_i+1)}{\alpha_i!}.
\]
These formulas are polynomial identities for real $x$, so allow
negative lattice offsets. The identity
$\binom{t+1}{a}-\binom ta=\binom t{a-1}$ implies
\[
 \Delta_i P_{j+1}(v)=P_j(\Delta_i v),\qquad
 P_j(v)(0)=0\ (j>0).
\]
They are also the polynomials obtained by the recursion
\[
 P_{j+1}(v)(x)=(j+1)^{-1}\sum_i x_iP_j(\Delta_i v)(x-e_i).
\]
To check this, expand the falling factorials and use $\sum_i\alpha_i=j+1$.
Their coefficients give
\[
 \|P_j(x)\|\le c_jM\theta L^{-j-1}(1+|x|_1)^j,
                         \qquad c_j=[4(j+1)]^j.
\]
Indeed each falling-factorial factor is bounded by
$(1+|x|_1+j)^j$, while summing the reciprocal factorials at
fixed degree contributes $4^j/j!\le4^j$.

If $|x|_1\le L$, all differences of order $d$ in the ball of
radius $L$ are at most $2^pM\theta L^{-d-1}$. The remainder
$R_d=v-\sum_{j<d}P_j$ is zero at zero for $d>0$, and
$\Delta_iR_d(v)=R_{d-1}(\Delta_i v)$. Induction on $d$,
followed by summation along a shortest signed coordinate path
from zero to $x$, gives
\[
 \|R_d(x)\|\le2^pM\theta L^{-d-1}(1+|x|_1)^d.
\]
For a negative step the forward difference is evaluated at the
new endpoint, which still lies in the same ball. Thus this
summation does not assume positive offsets. Since $p\ge d$,
it is bounded by the remainder in \eqref{eq:moment-newton}.
If $|x|_1>L$, then
\[
 (1+|x|_1/L)^p\le2^p L^{-p}(1+|x|_1)^p
\]
and, for $j<d$,
$L^{d-j}\le(1+|x|_1)^{d-j}$. Apply these inequalities to
\eqref{eq:moment-drift} with the empty word and to the piece
bounds, then use $R_d=v-\sum_{j<d}P_j$. This gives
\eqref{eq:moment-newton} with
$C_{d,p}=2^p+\sum_{j<d}c_j$. When $d=0$ only the remainder
is present, and the original bound suffices.

Evaluate these pieces at the actual centered lift on the torus.
The lift satisfies $|\operatorname{lift}(y-o)|_1\le
4\operatorname{dist}(y,o)$. Multiplying the constant by $4^p$
therefore gives a global decomposition into $d$ polynomial
pieces and a last remainder. The screened $j$th piece belongs
to $\mathcal M_{j+1}$. Nothing is required of the polynomial
behavior of the last remainder, including across the lift boundary.

\subsection{Applying the decomposition to the actual kernel}
The coarse-to-fine linear field is the installed $A_0$ kernel,
with the physical normalization of Section~\ref{sec:kernel-transport-identities}.
Its iterated spatial differences obey the strip estimates
\[
 \|\Delta_w A_0\|_\alpha\le A L^{-1-|w|},
                                      \qquad |w|\le d.
\]
These are the $A_0$ difference bounds in
Section~\ref{sec:curvature-horizon-producer}; no such assertion
about all the other tower blocks is needed here. If a coarse
input obeys $\|u(y)\|\le\theta(1+\operatorname{dist}(y,o_c))^p$,
where $o_c$ is the parent of the chosen fine anchor $o$, use the
installed row at $o+\operatorname{wrap}(x)$. Its summation addresses
are the parent of this site plus $\operatorname{wrap}(z)$, where
$z$ is the cover displacement in that residue row. The parent
map can cross a block boundary; its centered residue changes at
the same time. The coordinate division with remainder gives
\[
 1+\operatorname{dist}\bigl(\operatorname{parent}
 (o+\operatorname{wrap}(x))+\operatorname{wrap}(z),o_c\bigr)
 \le3(1+|x|_1/L)(1+|z|_1).
\]
Indeed in each coordinate the parent displacement is the integer
quotient of $x$ plus a residue correction of absolute value at
most one; periodic distance can only decrease this bound. Taking
the maximum of the coordinate bounds and then adding the cover
shift gives the displayed factor three. Taking absolute values
in the convolution now costs $3^p(1+|x|_1/L)^p$ times the strip
row sum with the additional weight $(1+|z|_1)^p$. The inequality
$(1+t)^p\le C_{p,\alpha}e^{\alpha t}$ absorbs that weight into
the available strip estimate. Sum also over the fixed form types.
Differences can
be taken on the row before summation by absolute convergence;
the installed odd-block carry identifies them with the cover
forward differences. We obtain \eqref{eq:moment-drift}, with
$M$ depending only on $A,\alpha,p$ and the fixed component count.
The preceding decomposition therefore applies to this actual
field, uniformly in the lattice size and block factor.

\subsection{Quantitative cancellation and displacement}
Let $\|\cdot\|_P$ denote the direction norm
\[
 \|v\|_P=\sup_e
 \frac{\|v_e\|}{\theta(1+\operatorname{dist}(o,e_-))^p}.
\]
Say that a direction admits a graded decomposition at scale
$0\le\xi\le1$ and cost $C\ge1$ if
\[
 v=\sum_{a=1}^{d+1}v_a,\qquad
 \|v_a\|_P\le C\xi^a,\qquad
                   v_a|_X\in\mathcal M_a\quad(a\le d).
\]
The preceding construction gives this property with $\xi=L^{-1}$.
Suppose the regular derivative bound of $F=f\circ\exp$ is
\begin{equation}\label{eq:moment-regular-bound}
 |D^mF(b)[v_1,\ldots,v_m]|
           \le Q\tau\prod_i\|v_i\|_P,
                                      \qquad m\le n,
\end{equation}
on the logarithmic region used below, with $n\ge d+1$.
Here $Q$ is the fixed conversion constant from the local
regularity norm, and $\tau$ is its price. The chart and health
conditions in that norm are retained in
\eqref{eq:moment-regular-bound}.

At zero, expand each argument into its $d+1$ pieces. All terms
of total weight at most $d$ vanish by
\eqref{eq:polynomial-moment-null}. Every remaining term has a
factor at most $C^m\xi^{d+1}$. There are $(d+1)^m$ terms.
For $m\le n$ this proves
\[
 |D^mF(0)[v_1,\ldots,v_m]|
             \le Q\tau((d+1)C)^n\xi^{d+1}=:S.
\]
Now let the displacement $b$ have the same decomposition, and
suppose that the segment $tb$, $0\le t\le1$, remains in the
healthy logarithmic region, with the tested links and direction
widths bounded by one as required by the regular norm. At order
$n$, \eqref{eq:moment-regular-bound} bounds the derivative on
this segment by $S$, since each argument has norm at most
$(d+1)C\xi$ and $n\ge d+1$. If the order-$m+1$ derivative is
bounded there by $(n-m)S$, the fundamental theorem of calculus
applied to $t\mapsto D^mF(tb)[v_1,\ldots,v_m]$ bounds its
change from zero by $(n-m)S$. Its value at zero is at most $S$.
Descending induction gives the uniform bound $(n+1)S$ at every
order through $n$. This proves the attenuation after displacement;
it does not require differentiating the health indicator.
